\section{Introduction}
\label{sec:introduction}

Caputo derivatives of order $0<\alpha<1$ model hereditary effects in viscoelasticity, anomalous transport, and population dynamics~\cite{F01,F02,F03,F04}; fractional predator--prey models in particular form a sizeable literature~\cite{E06,E07,E08,E09}. An empirical difficulty sits underneath all of it. Estimating $\alpha<1$ from data does not establish a fractional mechanism, because on a bounded window $[0,T]$ a delay term $\xi(t-\tau)$ or a low-dimensional latent state $q\in\mathbb{R}^m$ generates trajectories within the noise floor of the fractional ones (Corollary~\ref{cor:closure} makes this quantitative). Design methods for fractional-order \emph{estimation} exist~\cite{FOED01,FOED02,FOED03,FOED04}; safety-constrained information acquisition has been analyzed in general settings~\cite{SAFE26A,SAFE26B}. The intersection --- discrimination between memory \emph{mechanisms}, under an ecological safety constraint --- is the subject here:

\begin{quote}
\emph{Which perturbations and measurements can distinguish fractional memory from delayed or latent alternatives in a strong-Allee ecological system, and when does ecological safety make that discrimination fundamentally difficult?}
\end{quote}

Safety is not decorative in this question. Below the Allee threshold $A$ the prey growth term $P(x)=rx(1-x/K)(x/A-1)$ turns negative and collapse is irreversible, so the perturbations most informative about memory ($\|u\|_\infty$ near the budget, broadband spectra) are the ones most likely to destroy the object of study. Section~\ref{sec:benchmark} measures this numerically: among the six classical excitation families, the best-discriminating input crosses the threshold in $100\%$ of stable-regime cells. Section~\ref{subsec:safe-search} then shows that this apparent conflict is a property of those families rather than of the ecology.



\paragraph{Provenance of the ecological backbone.}
The ecological backbone follows the companion certification study~\cite{P01}: the equations, locked parameters, coexistence equilibrium and Jacobian are restated here for self-containment and are not claimed as contributions of the present work. That study certifies dynamical regimes of a specified Caputo model; it does not address whether such a model can be told apart from a delayed or latent mechanism reproducing the same observations to within $\sigma$, which is the question posed here. No result below recertifies the backbone. With the ecology frozen, differences between simulated observations are attributable to the memory law or to the experiment, never to an altered population model. Only the integer-order and Caputo mechanisms share $J(A)$ exactly; the delayed and latent rivals are calibrated alternative response laws matched at the same operating point, input channel and units.

\paragraph{Main contributions.}
The paper is organised around one falsification and one positive design result, with a separate obstruction that survives both.

\begin{enumerate}
    \item \textbf{Safe waveforms outside the classical families outperform every classical baseline} (Section~\ref{subsec:safe-search}). Under the same peak-amplitude budget $\|u\|_\infty\le0.100$ and a hard Allee-margin constraint, a constrained search over $163\,840$ candidate waveforms returns piecewise-constant designs with four-class BIC macro-accuracy $0.803$ and zero observed threshold crossings, against $0.514$ for the only classical family that is safe at this budget and $0.763$ for the best classical design of any kind, which crosses in $35.7\%$ of cells. The gain is concentrated where the safe classical design fails: Caputo recall rises from $0.404$ to $0.911$, delay recall from $0.445$ to $0.853$.
    \item \textbf{The six-family safety--informativeness frontier is not a property of the system} (Sections~\ref{subsec:amplitude}--\ref{subsec:safe-search}). Reducing amplitude does not restore safety --- the leading design still crosses in $71.4\%$ of cells at half amplitude, while its realised response gain rises to $4.86$ times the linear certificate, the signature of basin escape rather than linear response. The severe frontier reported for classical excitation families is therefore a statement about that parameterisation. A Pareto cost remains: the best unsafe candidate separates the hypotheses substantially better than the best safe one, so safety is not free.
    \item \textbf{A separate finite-horizon latent-complexity obstruction, quantified at response level to $m=128$} (Sections~\ref{subsec:semantics}--\ref{subsec:state}). Computed directly on the prey response rather than inferred from the memory kernel, the $L^1$ surrogate error falls to $4.8\cdot10^{-5}$ at $m=128$ and the induced exact minimax testing floor reaches $0.49999$: at attainable latent orders the finite-state rival is indistinguishable under the declared protocol, for every input in the budget. Waveform search does not remove this. In the benchmark the leading safe design also recovers only $0.475$ recall against the finite-latent rival, which is consistent with the obstruction but not evidence for it: the BIC penalty on the higher-dimensional model is an alternative explanation that this study does not separate.
    \item \textbf{Trajectory safety status, stated exactly} (Sections~\ref{sec:design} and~\ref{sec:benchmark}). Safety of the reported designs is verified a posteriori by mesh refinement together with a rigorous inter-node modulus; for four of the eight headline trajectories a validated enclosure with self-consistent error control is obtained, and only those are called rigorous. Constants entering the nonlinear invariant-set route are interval-certified, and that route is shown to admit inputs $70$ times smaller than the amplitude actually used.
    \item \textbf{Reproducibility control.} The frozen four-class benchmark reproduces to four decimal places on an independent host under a re-implemented driver ($2268$ cells across three amplitudes, no solver divergences).
\end{enumerate}

Contributions 1, 2 and 5 are empirical, established by the benchmark and the search; contribution 3 combines a proved bound (Theorem~\ref{thm:T20}) with computed surrogate errors, interval-enclosed for $m\le32$ and high-accuracy numerical for $m=64,128$; contribution 4 is a verification result whose exact status is stated per trajectory in Section~\ref{sec:benchmark}.

The structural separation of Section~\ref{sec:analysis} and the positive exponential-sum approximation it is contrasted with are \emph{not} contributions: the former is classical fractional-systems material, and the latter is the construction of McLean~\cite{K07} building on Beylkin and Monz\'on~\cite{K03,K04}. We record only an endpoint-inclusive $L^1(0,T)$ variant with the exact horizon scaling the convolution argument requires (Lemma~\ref{lem:soe}), and attribute the technique where it is used.

Complete proofs, secondary results, the Bayesian sequential extension, and solver derivations are in the Supplementary Material; the main text keeps the four statements needed to read the numerics.

\begin{figure}[htbp]
\centering
\includegraphics[width=0.98\textwidth]{fig17_paper_workflow}
\caption{Structure of the paper. Sections~\ref{sec:model}--\ref{sec:analysis}: ecological model, rival memory laws, and the four analytical results. Sections~\ref{sec:numerics}--\ref{sec:response}: solver validation (Table~\ref{tab:solver-validation}) and frequency/finite-horizon evidence. Section~\ref{sec:fd}: fractional-delay simulations. Sections~\ref{sec:design}--\ref{sec:benchmark}: input design, safety, and the 1512-cell benchmark. Proofs and secondary theory: Supplementary Material.}
\label{fig:workflow}
\end{figure}
